\begin{center}\small
\begin{tabular}{rrrrrl}
\toprule
Variables & Degree & Residuals & Live operations & Height bound & Seed\\\midrule
194 & 6 & 131 & 2,682 & $2^{292}$ & $P_{211}$\\
147 & 8 & 84 & 2,559 & $2^{325}$ & $P_{211}$\\
135 & 12 & 72 & 2,498 & $2^{331}$ & $P_{211}$\\
121 & 16 & 58 & 2,453 & $2^{324}$ & $P_{211}$\\
115 & 24 & 52 & 2,438 & $2^{306}$ & $P_{209}$\\
111 & 32 & 48 & 2,429 & $2^{364}$ & $P_{209}$\\
105 & 64 & 42 & 2,495 & $2^{546}$ & $P_{207}$\\
97 & 256 & 34 & 2,398 & $2^{1670}$ & $P_{207}$\\
90 & 32,768 & 27 & 2,306 & $2^{114689}$ & $P_{206}$\\
\bottomrule\end{tabular}\end{center}
The height column is a checked upper bound on the maximum absolute expanded coefficient. 
The largest literal integer in each circuit is 384. Live operations count binary 
addition, subtraction, and multiplication in the emitted shared circuit, including squares.

The 90-variable equation has 27 residuals and exact degree 32,768. 
It is the smallest-arity equation produced in this package. The 121-variable 
degree-16 equation is a useful alternative when degree matters. Their exact files are 
\code{higher_degree/equations/integer90_degree32768.json} and 
\code{higher_degree/equations/integer121_degree16.json}.
